\documentclass[10pt,aspectratio=169]{beamer}
\input{../../_en_preambulo_beamer}% Comandos y macros estadísticas compartidas para las presentaciones Beamer en inglés (Metropolis)

% Conjuntos numéricos básicos
\newcommand{\R}{\mathbb{R}}
\newcommand{\N}{\mathbb{N}}
\newcommand{\Q}{\mathbb{Q}}
\newcommand{\Z}{\mathbb{Z}}
\newcommand{\set}[1]{\left\{ #1 \right\}}
\newcommand{\abs}[1]{\left|#1\right|}
\newcommand{\norm}[1]{\left\| #1 \right\|}

% Comandos estadísticos y probabilísticos del curso
\newcommand{\Var}{\operatorname{Var}}
\newcommand{\cov}{\operatorname{Cov}}
\newcommand{\corr}{\rho}
\newcommand{\comb}[2]{\begin{pmatrix} #1 \\ #2 \end{pmatrix}}
\newcommand{\s}{\sigma}
\newcommand{\particion}{\mathfrak{P}}
\newcommand{\card}[1]{\#\left( #1 \right)}
\newcommand{\seq}[3]{\set{#1}_{#2}^{#3}}
\newcommand{\E}{\mathbb{E}}
\newcommand{\Prob}{\mathbb{P}}

% Entornos pedagógicos y cajas en inglés académico natural (soportando tanto nombres en inglés como en español)
\newenvironment{discussion}{%
  \begin{alertblock}{Discussion Question}%
}{%
  \end{alertblock}%
}
\newenvironment{discusion}{%
  \begin{alertblock}{Discussion Question}%
}{%
  \end{alertblock}%
}

\newenvironment{reflection}{%
  \begin{alertblock}{Classroom Reflection}%
}{%
  \end{alertblock}%
}
\newenvironment{reflexion}{%
  \begin{alertblock}{Classroom Reflection}%
}{%
  \end{alertblock}%
}

\newenvironment{paradox}{%
  \begin{alertblock}{Exploratory Paradox}%
}{%
  \end{alertblock}%
}
\newenvironment{paradoja}{%
  \begin{alertblock}{Exploratory Paradox}%
}{%
  \end{alertblock}%
}

\newenvironment{motivation}{%
  \begin{exampleblock}{Motivation: Data Science in Practice}%
}{%
  \end{exampleblock}%
}
\newenvironment{motivacion}{%
  \begin{exampleblock}{Motivation: Data Science in Practice}%
}{%
  \end{exampleblock}%
}

\newenvironment{quickchallenge}{%
  \begin{block}{Quick Team Challenge}%
}{%
  \end{block}%
}
\newenvironment{retoclase}{%
  \begin{block}{Quick Team Challenge}%
}{%
  \end{block}%
}

\title[Variance tests]{Section 07.08: Variance Tests}\subtitle{Dispersion and homoscedasticity}
\author[J. Castillo Colmenares]{Juliho Castillo Colmenares}\institute{Tecnológico de Monterrey}\date{}
\begin{document}
\begin{frame}[plain]\titlepage\end{frame}
\begin{frame}{Roadmap}\small\begin{enumerate}\item One variance.\item Comparing two variances.\item Assumptions and interpretation.\end{enumerate}\end{frame}
\begin{frame}{Source and scope}\small This section presents dispersion tests for uniformity, risk, and homoscedasticity assumptions.\begin{block}{Guiding question}Does observed variability match the hypothesized variability?\end{block}\end{frame}
\begin{frame}{Population variance}\small For a normal sample of size $n$, the sample variance $S^2$ summarizes dispersion around $\bar X$.\end{frame}
\begin{frame}{One-variance hypotheses}\small\[H_0:\sigma^2=\sigma_0^2\qquad\text{versus an alternative about }\sigma^2.\]\end{frame}
\begin{frame}{Chi-square statistic}\small\[\chi^2=\frac{(n-1)S^2}{\sigma_0^2}.\]\begin{block}{Distribution} Under $H_0$, $\chi^2\sim\chi^2_{n-1}$.\end{block}\end{frame}
\begin{frame}{Rejection regions}\small The alternative determines a right-tail, left-tail, or two-tail region of the chi-square distribution.\end{frame}
\begin{frame}{Two variances}\small For two independent normal samples, test $H_0:\sigma_1^2=\sigma_2^2$ using the ratio of sample variances.\end{frame}
\begin{frame}{$F$ statistic}\small\[F=\frac{S_1^2}{S_2^2}.\]\begin{block}{Distribution} Under equality, $F\sim F_{n_1-1,n_2-1}$.\end{block}\end{frame}
\begin{frame}{Ratio order}\small Put the larger variance in the numerator for a convenient right-tail test; keep degrees of freedom in the same order.\end{frame}
\begin{frame}{Procedure I}\small\begin{enumerate}\item Check normality and independence.\item State hypotheses.\item Set $\alpha$.\end{enumerate}\end{frame}
\begin{frame}{Procedure II}\small\begin{enumerate}\setcounter{enumi}{3}\item Compute $S^2$ or $F$.\item Obtain a quantile or $p$-value.\item Conclude in context.\end{enumerate}\end{frame}
\begin{frame}{Assumptions}\small $\chi^2$ and $F$ tests are sensitive to non-normality. Independence and the sampling design are essential.\end{frame}
\begin{frame}{Application: one variance}\small Given $\sigma_0^2$, $n$, and $S^2$, compute $\chi^2=(n-1)S^2/\sigma_0^2$.\end{frame}
\begin{frame}{Application: decision}\small Compare the statistic with $\chi^2_{n-1}$ quantiles according to the alternative.\end{frame}
\begin{frame}{Application: two variances}\small Two processes yield $S_1^2$ and $S_2^2$. The ratio $F$ evaluates whether their dispersions can be treated as equal.\end{frame}
\begin{frame}{Application: interpretation}\small Rejecting equal variances indicates different dispersion; it does not identify the cause or process mechanism.\end{frame}
\begin{frame}{Homoscedasticity}\small Equal variances support comparable errors across groups in parametric models. The test is diagnostic, not a replacement for graphical analysis.\end{frame}
\begin{frame}{Common errors}\small\begin{itemize}\item Omit normality.\item Reverse degrees of freedom when reordering $F$.\item Interpret dispersion as a mean difference.\end{itemize}\end{frame}
\begin{frame}{Reporting template}\small\begin{block}{Report} ``We obtained $\chi^2/F=\cdots$, with $df=\cdots$ and $p=\cdots$. At $\alpha=\cdots$, we [reject/do not reject] the stated equality.''\end{block}\end{frame}
\begin{frame}{Synthesis}\small\begin{itemize}\item One variance uses $\chi^2$.\item Two variances use $F$.\item Assumptions determine inferential validity.\end{itemize}\end{frame}
\begin{frame}{Chapter connection}\small These tests support homoscedasticity assessment in mean comparisons and regression models.\end{frame}
\end{document}
